\documentclass[12pt,a4paper]{article}

\usepackage[T1]{fontenc}
\usepackage[utf8]{inputenc}
\usepackage[romanian]{babel}
\usepackage{geometry}

\geometry{margin=2.5cm}

\title{Tutoriat 3\\Cautare Adversariala - Alpha-Beta Pruning}
\author{}
\date{}

\begin{document}

\maketitle

\section{Retezări în arborele de joc}

Pe măsură ce adâncimea arborelui crește, numărul de stări analizate de MiniMax devine foarte mare.

O idee importantă este eliminarea ramurilor care nu pot influența decizia finală.

Acest proces poartă numele de retezare (pruning).

\subsection{Idee}

În timpul explorării arborelui pot exista noduri care nu mai trebuie analizate deoarece rezultatul lor nu va modifica alegerea optimă.

Aceste ramuri sunt ignorate.

\subsection{Avantaje}

\begin{itemize}
\item reduce numărul de noduri analizate;
\item reduce timpul de execuție;
\item permite explorarea unei adâncimi mai mari;
\item păstrează aceeași soluție ca MiniMax.
\end{itemize}

\subsection{Exemplu}

Considerăm un nod MIN care nu a fost explorat complet dar are deja valoarea:

\[
2
\]

Dacă MAX are deja o alternativă cu valoarea:

\[
5
\]

atunci ramura respectivă nu mai poate influența decizia finală și poate fi eliminată din analiză.

\subsection{Concluzie}

Retezarea permite reducerea costului algoritmului MiniMax prin evitarea explorării ramurilor inutile.

Cea mai cunoscută tehnică de retezare este algoritmul Alpha-Beta Pruning.







\section{Algoritmul Alpha-Beta Pruning}

Alpha-Beta Pruning reprezintă o optimizare a algoritmului MiniMax care reduce numărul de noduri analizate fără a modifica rezultatul final.

Algoritmul obține aceeași mutare optimă ca MiniMax, dar într-un timp mai scurt.

\subsection{Ideea principală}

În timpul explorării arborelui, anumite ramuri pot fi eliminate deoarece nu mai pot influența decizia finală.

Aceste ramuri nu mai sunt evaluate.

\subsection{Parametrii Alpha și Beta}

Algoritmul utilizează două valori:

\begin{itemize}
\item $\alpha$ - cea mai bună valoare găsită până acum pentru MAX;
\item $\beta$ - cea mai bună valoare găsită până acum pentru MIN.
\end{itemize}

Pe parcursul explorării:

\[
\alpha \leq \beta
\]

Dacă această condiție nu mai este respectată, căutarea se oprește pe ramura respectivă.

\subsection{Condiția de retezare}

O ramură poate fi eliminată atunci când:

\[
\alpha \geq \beta
\]

În acel moment știm că explorarea continuă nu poate schimba decizia finală.

\subsection{Avantaje}

\begin{itemize}
\item reduce numărul de noduri analizate;
\item reduce timpul de execuție;
\item permite explorarea unor adâncimi mai mari;
\item produce exact aceeași soluție ca MiniMax.
\end{itemize}

\subsection{Dezavantaj}

Implementarea este puțin mai complexă decât cea a algoritmului MiniMax.

\subsection{Concluzie}

Alpha-Beta Pruning reprezintă cea mai importantă optimizare a algoritmului MiniMax și permite rezolvarea eficientă a jocurilor cu arbori de dimensiuni mari.














\section{Implementarea Alpha-Beta Pruning}

Alpha-Beta Pruning reprezintă o versiune optimizată a algoritmului MiniMax.

Diferența principală constă în utilizarea limitelor Alpha și Beta pentru eliminarea ramurilor care nu mai pot influența rezultatul final.

\subsection{Pseudocod}

\begin{verbatim}
AlphaBeta(stare, alpha, beta, jucator)

    daca stare este terminala
        returneaza utilitatea(stare)

    daca jucator este MAX

        valoare <- -infinit

        pentru fiecare succesor

            valoare <- max(
                valoare,
                AlphaBeta(
                    succesor,
                    alpha,
                    beta,
                    MIN
                )
            )

            alpha <- max(alpha, valoare)

            daca alpha >= beta
                opreste cautarea

        returneaza valoare

    altfel

        valoare <- +infinit

        pentru fiecare succesor

            valoare <- min(
                valoare,
                AlphaBeta(
                    succesor,
                    alpha,
                    beta,
                    MAX
                )
            )

            beta <- min(beta, valoare)

            daca alpha >= beta
                opreste cautarea

        returneaza valoare
\end{verbatim}

\subsection{Implementare în Python}

\begin{verbatim}
def alpha_beta(node,
               depth,
               alpha,
               beta,
               maximizing_player):

    if depth == 0 or node.is_terminal():
        return node.utility()

    if maximizing_player:

        value = float("-inf")

        for child in node.children:

            value = max(
                value,
                alpha_beta(
                    child,
                    depth - 1,
                    alpha,
                    beta,
                    False
                )
            )

            alpha = max(alpha, value)

            if alpha >= beta:
                break

        return value

    else:

        value = float("inf")

        for child in node.children:

            value = min(
                value,
                alpha_beta(
                    child,
                    depth - 1,
                    alpha,
                    beta,
                    True
                )
            )

            beta = min(beta, value)

            if alpha >= beta:
                break

        return value
\end{verbatim}

\subsection{Rolul parametrilor}

\begin{itemize}
\item alpha reprezintă cea mai bună valoare găsită până acum pentru MAX;
\item beta reprezintă cea mai bună valoare găsită până acum pentru MIN;
\item depth reprezintă limita de adâncime a căutării.
\end{itemize}

\subsection{Avantaj față de MiniMax}

Alpha-Beta Pruning produce exact aceeași soluție ca MiniMax, dar analizează mai puține noduri.

Prin eliminarea ramurilor inutile, algoritmul poate explora arbori mai adânci în același timp de execuție.

\subsection{Concluzie}

Alpha-Beta Pruning este cea mai utilizată optimizare a algoritmului MiniMax și reprezintă baza motoarelor moderne pentru jocuri precum șahul, damele sau X și 0.








\section{Intuiția algoritmului Alpha-Beta}

Alpha-Beta Pruning are aceeași idee de bază ca MiniMax, dar evită explorarea ramurilor care nu mai pot influența decizia finală.

\subsection{Intuiție}

În timpul căutării, dacă știm deja că o anumită mutare este mai slabă decât una analizată anterior, nu mai are sens să continuăm explorarea acelei ramuri.

Aceasta poate fi eliminată.

\subsection{Exemplu}

Considerăm arborele:

\begin{verbatim}
          MAX
         /   \
       MIN   MIN
      / \    / \
     3   5  2   ?
\end{verbatim}

Ramura stângă oferă:

\[
MIN(3,5)=3
\]

Prin urmare:

\[
\alpha = 3
\]

Când analizăm ramura dreaptă, observăm imediat valoarea:

\[
2
\]

Întrucât nodul este MIN, valoarea finală a ramurii drepte va fi cel mult 2.

MAX deține deja alternativa:

\[
3
\]

Astfel, restul ramurii poate fi ignorat deoarece nu va fi niciodată aleasă.

\subsection{Rezultat}

Alpha-Beta Pruning reduce numărul de noduri analizate fără a modifica alegerea finală.

Algoritmul obține exact aceeași soluție ca MiniMax, dar într-un timp mai scurt.

\subsection{Concluzie}

Ideea esențială este simplă:

\begin{quote}
Nu analizăm ramurile despre care știm deja că nu pot influența decizia finală.
\end{quote}

Această optimizare permite explorarea unor arbori mult mai mari decât în cazul MiniMax clasic.




\section{Exemple Alpha-Beta}

\subsection{Exemplul 1}

Considerăm arborele:

\begin{verbatim}
          MAX
         /   \
       MIN   MIN
      / \    / \
     3   5  2   9
\end{verbatim}

Ramura stângă este evaluată prima:

\[
MIN(3,5)=3
\]

Astfel:

\[
\alpha = 3
\]

În ramura dreaptă întâlnim valoarea:

\[
2
\]

Cum nodul aparține lui MIN și:

\[
2 < \alpha
\]

restul ramurii nu mai poate influența decizia lui MAX.

Valoarea 9 nu mai trebuie evaluată.

\subsection{Exemplul 2}

Considerăm arborele:

\begin{verbatim}
                 MAX
              /       \
            MIN       MIN
           /  \      /  |  \
          3    5    2   4   8
\end{verbatim}

Pentru primul nod MIN:

\[
MIN(3,5)=3
\]

Prin urmare:

\[
\alpha = 3
\]

La al doilea nod MIN este analizată mai întâi valoarea:

\[
2
\]

Deoarece:

\[
2 < 3
\]

știm că rezultatul acelui nod va fi cel mult 2.

Valorile 4 și 8 nu mai trebuie analizate.

\subsection{Comparație cu MiniMax}

MiniMax:

\begin{itemize}
\item analizează toate nodurile;
\item garantează soluția optimă.
\end{itemize}

Alpha-Beta:

\begin{itemize}
\item elimină ramurile inutile;
\item găsește aceeași soluție;
\item analizează mai puține noduri.
\end{itemize}

\subsection{Concluzie}

Pe măsură ce dimensiunea arborelui crește, numărul de noduri eliminate poate deveni foarte mare.

De aceea Alpha-Beta Pruning este mult mai eficient decât MiniMax pentru jocurile complexe.

















\section{Proprietățile algoritmului Alpha-Beta}

Alpha-Beta Pruning păstrează toate proprietățile algoritmului MiniMax, dar reduce numărul de noduri analizate.

\subsection{Completitudine}

Algoritmul este complet.

Dacă există o soluție și arborele este finit, aceasta va fi găsită.

\subsection{Optimalitate}

Algoritmul este optim.

Rezultatul obținut este identic cu cel furnizat de MiniMax.

Retezarea unor ramuri nu modifică mutarea finală aleasă.

\subsection{Complexitatea în timp}

În cel mai rău caz:

\[
O(b^d)
\]

adică aceeași complexitate ca MiniMax.

În practică însă, numărul de noduri analizate este mult mai mic.

În cel mai bun caz:

\[
O(b^{d/2})
\]

ceea ce permite explorarea unor arbori semnificativ mai adânci.

\subsection{Complexitatea în spațiu}

La fel ca MiniMax:

\[
O(b \cdot d)
\]

deoarece arborele este explorat folosind o strategie de tip Depth First Search.

\subsection{Avantaje}

\begin{itemize}
\item produce aceeași soluție ca MiniMax;
\item analizează mai puține noduri;
\item reduce timpul de execuție;
\item permite căutări mai adânci.
\end{itemize}

\subsection{Limitări}

\begin{itemize}
\item eficiența depinde de ordinea în care sunt analizate mutările;
\item în cel mai rău caz nu oferă îmbunătățiri față de MiniMax;
\item implementarea este puțin mai complexă.
\end{itemize}

\subsection{Concluzie}

Alpha-Beta Pruning reprezintă cea mai eficientă optimizare clasică a algoritmului MiniMax.

Acesta păstrează completitudinea și optimalitatea, reducând în același timp numărul de noduri care trebuie explorate.


\end{document}